\section{Subdivisions determined by actual corners}
\label{sec:area_law_actual_side_mask}

The opposing corners determine the subdivisions used in the cell fans;
see \cite[Section~11]{OpenAI2026AreaLaw}. Retain the common origin, endpoint
set and dyadic layers, and fix a lower index $k_0$. All actual fine cells
at indices at least $k_0$ participate in this construction.

% Source: 10-geometry.tex, lines 299--310. Dummy-neighborhood corners
% and opposing-region constancy are separate assertions.

\begin{theorem}[Geometry of an elementary side]
    \label{thm:area_law_elementary_side_geometry}
    \lean{TNLean.PEPS.AreaLaw.Geometry.cellFan_elementary_endpoints_cases,
        TNLean.PEPS.AreaLaw.Geometry.cellFan_elementary_segment_subset_whole}
    \leanok
    \uses{def:area_law_cell_fan}
    Fix any dyadic cell and optional midpoint subdivision. Let $[a,b]$
    be a whole side and $[u,v]$ one of its elementary segments, ordered
    counterclockwise. One of the following descriptions applies:
    the side is undivided and $(u,v)=(a,b)$; the side is divided and
    $(u,v)=(a,(a+b)/2)$; or the side is divided and
    $(u,v)=((a+b)/2,b)$. In every case $[u,v]\subseteq[a,b]$.
\end{theorem}
\begin{proof}\leanok
    \uses{def:area_law_cell_fan}
    In the affine parametrization of the whole side by $[0,1]$, the
    elementary endpoints have parameters $0,1$ for an undivided side,
    $0,1/2$ for the first half, and $1/2,1$ for the second half.
    Each resulting interval belongs to $[0,1]$.
\end{proof}

\begin{definition}[The actual midpoint subdivision]
    \label{def:area_law_actual_side_mask}
    \lean{TNLean.PEPS.AreaLaw.Geometry.fineLayerSplitMask,
        TNLean.PEPS.AreaLaw.Geometry.fineLayerSplitMask_eq_true_iff}
    \leanok
    \uses{def:area_law_cell_fan,def:area_law_cell_corner,
        def:area_law_fine_layer_indices}
    For the fine cell indexed by $z$ in layer $k$, write $[a_s,b_s]$
    for its whole side $s$. Divide this side at its midpoint precisely
    when that midpoint is a corner of some actual fine cell in a layer
    $h\ge k_0$. The elementary segments are the whole side if no such
    corner exists, and its two midpoint halves otherwise.
\end{definition}

\begin{theorem}[Opposing corners are elementary endpoints]
    \label{thm:area_law_elementary_corner_endpoint}
    \lean{TNLean.PEPS.AreaLaw.Geometry.fineLayer_corner_on_elementarySide}
    \leanok
    \uses{def:area_law_actual_side_mask}
    Suppose that $C_0\ge2$, $k\ge50\,000\,000$, and $z$ indexes an
    actual fine cell of layer $k$. Let $w$ index an actual fine cell of
    any layer $h\ge k_0$. If a corner of the cell indexed by $w$ belongs
    to an elementary segment of the reference cell under the actual
    midpoint subdivision, then it is an endpoint of that segment.
\end{theorem}
\begin{proof}\leanok
    \uses{thm:area_law_corner_side_subdivision,def:area_law_actual_side_mask,
        thm:area_law_elementary_side_geometry}
    Every elementary segment lies on its whole side. The whole-side
    assertion gives its start, midpoint or end as the only possible
    opposing corners. For an undivided side, a corner at the midpoint
    would be a witness requiring subdivision. For either midpoint half,
    the opposite endpoint of the whole side lies outside that half.
    The remaining two possibilities are exactly its endpoints.
\end{proof}
